Our methodology is designed to assess the security of SSH client signatures in the wild. As actual client signatures can only be observed by trusted servers, they can not be collected directly. We therefore approach this issue in multiple steps. First, we collect a large number of SSH public keys from Git-based online collaborative development platforms, namely GitHub, GitLab, and Launchpad, as described in \autoref{sec:methodology:sub:collection}. The fact that these platforms reveal the SSH public keys of their users has previously been noted by engineers~\cite{cox2015,golubin2019,filippo2015}, who highlighted the risk to user's privacy.\footnote{As a practical demonstration, Filippo Valsorda maintains a special SSH server \texttt{whoami.filippo.io}, that identifies users trying to connect to it by comparing their public key to a mirror of GitHub's user database.} We then subject these keys to an extensive test suite based on the rules of the CA/Browser forum for TLS server certificates~\cite{ca-browser-forum-2.1.2}, detailed in \autoref{subsec:securitykeys}. Compared to previous work, which limits its scope to specific vulnerabilities or key types, we provide the first comprehensive analysis of all SSH public keys on these platforms. In \autoref{sec:methodology:sub:restrictions}, we analyze these platforms for upload restrictions that are designed as a protection mechanism to prevent users from uploading weak or compromised keys.

Following the analysis of public keys, we also conduct extensive tests of SSH clients and agents in a controlled environment, as described in \autoref{sec:methodology:sub:clients}. In particular, we focus on identifying the underlying cryptographic library and supported algorithms for each analyzed SSH client, followed by checks for deterministic or biased nonces. 
%Together, this allows us to estimate the security of SSH client signatures in the wild.

\subsection{Large-Scale Collection of SSH Client Keys}
\label{sec:methodology:sub:collection}

For our analysis in \autoref{sec:clientkeys}, we collect publicly accessible SSH client keys from the online collaborative development platforms GitHub, GitLab, and Launchpad. These platforms commonly use SSH keys to authenticate users accessing source code repositories for reading and writing. We exclude Bitbucket from our study, although it is the third-largest platform, because downloading user SSH keys requires authorized access to workspaces. For each platform, keys are collected through the official GraphQL or REST APIs, respecting the platform's API limits. We use the GraphQL API for GitHub, and a combined approach for GitLab, using GraphQL to list users and REST to retrieve their SSH public keys. Launchpad only offers a REST API.

Each API response is stored in an Elasticsearch database for further processing. As API responses typically list SSH public keys by user, entries are transposed to list users by SSH public key instead. The resulting entries are checked for duplicate keys and, if no duplicate is found, stored in a separate index for analysis. In the case of duplicates, the existing entry is extended, which allows us to analyze public key reuse across multiple user accounts and platforms.

Using this methodology, we retrieve two raw data sets by scraping publicly available SSH keys from GitHub, GitLab, and Launchpad, the first in June 2023 and the second in January 2025 (\autoref{tab:longitudinal}).

\subsection{Security of SSH Client Keys}
\label{subsec:securitykeys}

We analyze each retrieved key according to the rules stated by the CA/Browser forum~\cite[Sections 6.1.1.3, 6.1.5, 6.1.6]{ca-browser-forum-2.1.2}. Although these rules are only intended for TLS server certificates, we believe long-lived SSH client key pairs should also fulfill those requirements related to the security of the key material. We also analyze the set of all RSA keys for common prime factors using the batch GCD attack. To the best of our knowledge, together, these tests cover all known, efficient attacks against public keys that can be run against a large database of keys. Other tests, for example, those found additionally in the RsaCtfTool,\footnote{\url{https://github.com/RsaCtfTool/RsaCtfTool}, accessed: 2025-04-14} are too slow to be run against a large dataset.

\paragraph{Short RSA keys}

\cite[Section 6.1.5]{ca-browser-forum-2.1.2} proposes that the encoded modulus length must be at least 2048 bits and fit byte boundaries. If the modulus $N=pq$ is too small, a brute-force attack that factors $N$ can recover an RSA private key from the public key. State-of-the-art is a number field sieve attack, such as CADO\nobreakdashes-NFS~\cite{cadonfs}. The largest RSA modulus currently factored is RSA\nobreakdashes-250 from the RSA factoring challenge~\cite{rsachallenge} and has 829 bits~\cite{boudot2020}. It was factored by Boudot et al. in 2020, using 2700 CPU-core-years. We consider all keys up to that size broken by brute-force factoring.

\paragraph{RSA Public Exponents}

\cite[Section 6.1.6]{ca-browser-forum-2.1.2} proposes that the public exponent $e$ must be an odd integer larger or equal to $3$ and should be in the range $[2^{16}+1,2^{256}-1]$. A key that is not compliant with this rule is suspect, because a large public exponent $e$ (i.e., roughly the same length as the modulus) may indicate that a small $d$ was chosen during key generation, and $e$ was computed using the extended Euclidean algorithm.
Since small values of $d$ can be found by exhaustive search, we check if the set of public RSA keys contains any large public exponent.

\paragraph{RSA Batch GCD and Small Prime Factors}

In 2012, Heninger et al.~\cite{USENIX:HDWH12} published an attack using a batch GCD algorithm to efficiently find prime factors of individual keys in large sets of RSA keys. The attack exploits cases where RSA keys share a common prime factor. We use an enhanced batch GCD algorithm incorporating subtrees to optimize memory usage, following an approach similar to that proposed by Hastings et al.~\cite{hastings2016}. Our implementation enables the subsequent addition of new keys and the testing of individual keys, with or without their inclusion in the dataset.

According to~\cite[Section 6.1.6]{ca-browser-forum-2.1.2}, the modulus $N$ should not have (prime) factors smaller than $752$. We extend this and check for the first one million prime factors ($p_\text{last}=15{,}485{,}867$) by adding their product during batch GCD computation. For all RSA moduli with one or multiple factors in the first ten thousand primes ($p_\text{last} = 104{,}743$), we apply the elliptic curve factorization algorithm by Lenstra~\cite{lenstra1987factoring} with conservative parameters suitable for finding factors up to 25 decimal digits, as previously done by Cryptosense~\cite{cryptosense2015}.

\paragraph{Debian, ROCA, and Fermat Vulnerabilities}

\cite[Section 6.1.1.3]{ca-browser-forum-2.1.2} states that keys should be checked for known weak keys, such as Debian weak keys~\cite{cve-2008-0166}, ROCA weak keys~\cite{CCS:NSSKM17}, and Fermat weak keys~\cite{EPRINT:Bock23}. In 2008, the OpenSSL version in Debian GNU/Linux had an inadequate procedure to initialize the random number generator. As a result, all RSA and (EC)DSA keys generated with such a system lack entropy and can be brute-forced with only a few hours of CPU time. The badkeys project\footnote{\url{https://badkeys.info/}, accessed: 2025-04-14} by Hanno Böck includes precomputed tables for 1024-bit DSA keys generated by OpenSSH, ECDSA keys under the NIST~\ptwofivesix and \pthreeeightfour curves generated by OpenSSL, as well as RSA keys of size 1024, 2048, 3072, and 4096-bit generated by OpenSSL and OpenSSH. We use code from the badkeys project to check for these keys while simultaneously testing for other known compromised keys. ROCA~\cite{CCS:NSSKM17} describes a type of keys whose prime numbers are generated with reduced entropy. ROCA keys can be detected efficiently with a negligible false positive rate. We use the tool from the original authors to test the keys in our database. The Fermat attack~\cite{EPRINT:Bock23} can find $p$ and $q$ by brute-force search around $\sqrt{N}$ if $p$ and $q$ are close. To achieve this, the attack tries to represent $N = pq$ as the difference of two squares $N = (a + b)(a - b) = a^2 - b^2$ for some $a$ and $b$. To do so, $a$ is initialized to the largest integer $\leq \sqrt{N}$. Then, in what is considered a round in the attack, it is checked whether $a^2 - N$ is a perfect square. If not, $a$ is incremented by $1$, and the attack continues with the next round. We apply the Fermat attack with 100 rounds.

\paragraph{(EC)DSA}

For DSA, we perform an explicit domain parameter validation with partial validation of the generator and a full public key validation as described in~\cite[Section 4.1]{nist-sp-800-89} and~\cite[Section 5.3.1]{nist-sp-800-89}, respectively. Any prime $p$ with fewer than 2048 bits is considered weak. A prime with fewer than 795 bits, the current record of discrete logarithm computation for safe primes~\cite{boudot2019}, is considered broken. Additionally, we check for compromised keys, such as Debian weak keys, using the badkeys tool.

For ECDSA, we first validate that NIST~\ptwofivesix, \pthreeeightfour, or \pfivetwoone is used as the underlying elliptic curve following~\cite[Section 6.1.5]{ca-browser-forum-2.1.2} and that the public key is a valid encoding of an elliptic curve public key. Then, we perform an ECC full public key validation routine described in~\cite[Section 5.6.2.3.3]{nist-sp-800-56a-3} per~\cite[Section 6.1.6]{ca-browser-forum-2.1.2}. This ensures that the public key point is not the identity element, that the coordinates are within valid bounds, that the public key is a point on the curve, and that it has the correct subgroup order. Additionally, we check for compromised keys, such as Debian weak keys, using the badkeys tool.


\paragraph{EdDSA}
For EdDSA public keys, to the best of our knowledge, there are no official security considerations or attacks to consider. Nevertheless, we validate the public key following the decoding process described in~\cite[Section 5.1.3]{rfc8032}, and by calculating the public key point's order. As the decoding process occurs on every signature validation, we expect that keys not passing this test will be unusable in practice. Additionally, we check for compromised keys, such as Debian weak keys, using the badkeys tool.

\subsection{Evaluating Public Key Upload Restrictions}
\label{sec:methodology:sub:restrictions}

To test for existing public key upload restrictions put forth by the platforms, we generate test keys that specifically fail our criteria described in \autoref{subsec:securitykeys}. Each test key should violate a single test criterion to allow for fine-grained detection. Note, however, that this cannot be achieved for all requirements; for example, a key with invalid ECDSA encoding cannot pass any subsequent test. Each test key is then uploaded to the platforms, and the response is recorded. A successful upload is indicated by the platform listing the public key under the user's public keys; a successful authentication is not required.

\subsection{Evaluating SSH Clients in the Lab}
\label{sec:methodology:sub:clients}

We restrict the evaluation of client capabilities to the supported signature algorithms because they determine the strength of the client's authentication.

\paragraph{Selection of Clients and Agents}

We compile a list of SSH clients for evaluation (\autoref{tab:clients}) from a comparison of SSH implementations assembled by Max Horn,\footnote{\url{https://ssh-comparison.quendi.de/}, accessed: 2025-04-14} a Wikipedia page on SSH clients,\footnote{\url{https://wikipedia.org/wiki/Comparison_of_SSH_clients}, accessed: 2025-04-14} and implementations well-known to us. We exclude clients that are no longer actively maintained, as well as clients supporting neither ECDSA nor EdDSA. Evaluation is limited to each client's most recent version, and \nrClients clients in total; evaluating historical versions is out of scope.

For SSH agents, no comparative lists exist. The list in \autoref{tab:agents} is based on internet search results, and includes all agents mentioned on the Wikipedia page on SSH agents.\footnote{\url{https://wikipedia.org/wiki/Ssh-agent}, accessed: 2025-04-14}

\paragraph{Determining Client Features}

For each SSH client implementation, we gather the supported signature algorithms based on the documentation. We also identify the underlying cryptographic library using basic binary analysis, process monitoring, and publicly available information on the vendor's website. For binary analysis under Linux, we use a combination of the tools \texttt{ldd} for identifying required shared libraries, \texttt{nm} for listing static and dynamic symbols, and \texttt{strings} for searching character sequences. Similarly, for binary analysis under Windows, we utilize the \texttt{dumpbin} and \texttt{strings} tools provided by Microsoft. Our findings are detailed in the additional material in the artifacts.

\label{sec:detnoncetest}
\paragraph{Tests for Deterministic/Biased Nonces}

This evaluation is limited to the DSA, ECDSA, and EdDSA algorithms due to the use of a nonce, where even a slight bias in nonce generation can allow for efficient recovery of the secret key. While major SSH implementations have disabled DSA support by default, and the OpenSSH project intends to remove DSA support from its source code entirely in early 2025, some clients still support it for compatibility reasons.

We implement a novel tool in Go that determines the nonce generation method and potential nonce bias, even for proprietary clients without available source code. To do so, the tool captures DSA, ECDSA, or EdDSA client authentication signatures by starting a custom SSH server. If supported by the SSH client, the client is automatically invoked to connect to the server. Otherwise, the connection is set up manually or by using UI automation tools, such as an automatic mouse clicker. We have access to the client's private key, as this tool is designed for in-lab analysis. With the help of this private key, the nonce values are recovered from the sampled signatures and analyzed.

To check for known deterministic nonce schemes from RFC~6979 and \putty \puttyvulnversion ($k_\text{proto}$), our tool generates two nonces according to these schemes and compares them with the nonce value from the received signature. To check for random nonces, the tool enforces further authentication attempts within the same connection, using the same client public key, by indicating a partial authentication success. Partial success is part of the SSH standard and allows for in-protocol, multifactor authentication by sending \msguserauthfailure with a special boolean flag set. Depending on their internal logic, a client may try to authenticate twice or more using the same public key in the same session. Such a client computes the signature twice over the same payload. If these two signatures are identical, we have detected a new, unknown deterministic nonce scheme. Otherwise, we can conclude that the nonce generation is at least partially random.

Our tool can also test for biased nonce generation, applying two statistical tests: a single-bit Z\nobreakdashes-test to detect any single-bit bias, and a more sophisticated Rayleigh test based on an implementation by Daniel Bleichenbacher in Project Wycheproof \footnote{\url{https://github.com/C2SP/wycheproof}, accessed: 2025-04-14} to detect a possible multi-bit bias.
When collecting nonces for this analysis, we do not indicate partial success to avoid introducing bias by repeatedly sampling the same nonce, and instead sample each signature over a fresh payload. The significance level of the test is set to $\alpha = 2^{-32}$ to avoid false positives. For each implementation, a minimum of 10,000 nonces is sampled for bias analysis.
